\section{Related work}
\label{sec:related}

\paragraph{Return-conditioned sequence modeling.}
Decision Transformer casts reinforcement learning as conditional sequence modeling over return-to-go, states, and actions \citep{chen2021decision}. Trajectory Transformer similarly studies offline RL through sequence modeling and planning over learned trajectory distributions \citep{janner2021trajectory}. Our study uses the return-conditioned inference interface from this line of work, but intentionally strips away benchmark complexity to isolate one selection mechanism.

\paragraph{Offline RL and support constraints.}
Offline RL methods address the danger of evaluating or improving policies outside the support of the behavior data \citep{levine2020offline}. D4RL provides benchmark datasets for this setting \citep{fu2021d4rl}. Conservative and support-aware methods such as BCQ, BEAR, CQL, and MOPO constrain or penalize out-of-distribution actions and values \citep{fujimoto2019bcq,kumar2019bear,kumar2020cql,yu2020mopo}. Our repair ladder is not a new offline RL algorithm; it borrows the support-awareness principle to audit return-conditioned trajectory selection.

\paragraph{Score reranking, reward models, and proxy overoptimization.}
Score-based reranking is widely used with learned reward models in language-model alignment and summarization systems \citep{stiennon2022summarize,nakano2022webgpt,ouyang2022instructgpt}. Recent analyses study theoretical properties of candidate reranking policies and the scaling of reward-model overoptimization \citep{beirami2025bon,gao2022rewardoveropt}. Our contribution is complementary: we give finite, tie-aware utility accounting for a fixed candidate distribution and apply it to return-conditioned trajectory selection with explicit offline-support diagnostics.

\paragraph{Calibration and tail risk.}
Calibration work emphasizes that high model confidence or score does not automatically imply correctness under distribution shift \citep{guo2017calibration}. Order statistics and extreme-value theory formalize how maxima concentrate in tails \citep{david2003order,coles2001extreme}. We use this perspective to explain why increasing $N$ is a tail-pressure intervention, and why the sign of utility change depends on the relation between proxy score, real utility, and support in that tail.

\paragraph{Specification and reward misspecification.}
The broader safety literature has long noted that optimized proxies can diverge from intended objectives \citep{amodei2016safety}. Our paper gives a small, reproducible instance of that point for return-conditioned policies: the optimized quantity is selected proxy return, while the intended quantity is real utility.
